\begin{abstract}
This paper proposes a fractional stochastic transportation model that integrates long-memory dynamics, demand uncertainty, and carbon trading within a unified sustainable transportation planning framework. The transportation system is modeled through a Caputo fractional balance equation, enabling the representation of persistent operational effects such as congestion carry-over, delayed recovery, and cumulative deterioration. Demand uncertainty is captured through a finite-scenario stochastic process, while environmental regulation is incorporated using a cap-and-trade mechanism. The resulting optimization problem is formulated as a two-stage mixed-integer stochastic program with Conditional Value-at-Risk (CVaR) to address risk aversion. A deterministic equivalent of the fractional model is derived using the L1 discretization scheme, and theoretical properties, including existence of optimal solutions, no-simultaneous carbon trading, and scenario stability, are established. To efficiently solve the resulting large-scale problem, a Sample Average Approximation combined with Progressive Hedging is developed. Computational experiments demonstrate that fractional memory significantly influences transportation decisions. Stronger memory leads to more conservative shipment policies, reducing emergency transportation and carbon emissions while increasing inventory-related costs. Sensitivity analyses further show that higher demand volatility increases both expected cost and risk, whereas tighter carbon allowances encourage endogenous allowance trading and operational adaptation. Comparisons with the classical first-order model confirm that incorporating long-memory effects provides a more realistic representation of transportation systems operating under uncertainty and environmental constraints. These results highlight the value of fractional modeling for sustainable and risk-aware transportation planning.
\end{abstract}

\noindent\textbf{Keywords:} fractional transportation, Caputo derivative, stochastic demand, carbon trading, CVaR, sample-average approximation, progressive hedging, sustainable logistics, long-memory systems.
